\documentclass[12pt]{article}
\usepackage[utf8]{inputenc}
\usepackage[T1]{fontenc}
\usepackage{hyperref}
\title{Uri Language Resources for K-12 Curriculum Development}
\author{Compiled by Claude Sonnet 5 (claude-sonnet-5) for Andr\'as Kornai}
\date{}
\begin{document}
\maketitle

\section{Introduction}
Uri (ISO 639-3 uvh; also called Erap, Uri Vehees, or Urii) is a Finisterre language of the Trans-New Guinea family, spoken in Morobe Province, Papua New Guinea. Wikipedia cites 2,500 native speakers from a 1991 count, and another source gives about 3,000 in the Boana district. This dossier is short because few resources turned up.

\section{Primary Educational Resources}
No primers or school materials were found.

\section{Linguistic and Reference Materials}
\begin{itemize}
\item Wikipedia overview: \url{https://en.wikipedia.org/wiki/Uri_language}
\item Numerals page, Max Planck Institute: \url{https://lingweb.eva.mpg.de/channumerals/Uri.htm}
\item ASJP word list: \url{https://asjp.clld.org/languages/URI_PAPUA_NEW_GUINEA}
\item Trans New Guinea database (Uri word entries): \url{https://transnewguinea.org/word/sand?page=13}
\item Research Data Australia record, Uri (Erap) language and culture: \url{https://researchdata.edu.au/uri-erap-language-culture/2045963}
\end{itemize}

\section{Religious and Community Resources}
The PARADISEC collection TW1 of the SIL linguist Tom Webb (active from about 1964 to 2000) contains manuscripts, audio recordings, and a New Testament translation in Uri. It is closed access: \url{https://catalog.paradisec.org.au/repository/TW1}

\section{Digital and Audio Resources}
Wikitongues language page: \url{https://wikitongues.org/languages/uvh/}

\section{Curriculum Development Context}
Uri is not known to be taught in schools. The most substantial material is the closed-access PARADISEC collection, so any curriculum work would need to negotiate access with the depositor and the community.
\end{document}